\documentclass[10pt,a4paper]{amsart}
\usepackage[margin=19mm]{geometry}
\usepackage{amsmath,amssymb,mathtools}
\usepackage[scr=rsfs,cal=euler]{mathalfa}
\usepackage{xfrac}
\usepackage[unicode,hidelinks,bookmarks=false]{hyperref}
\newcommand{\ie}{i.e.\ }
\newcommand{\cf}{cf.\ }
\newcommand{\resp}{respectively\ }
\newcommand{\Subsection}{\subsection}
\providecommand{\nobreakdash}{\nobreak\hbox{-}}
\setlength{\emergencystretch}{2em}
\multlinegap0pt
\usepackage{amssymb}
\newtheorem{lemma}{Lemma}
\newtheorem{fact}{Fact}
\title{Constructing edge-disjoint Steiner trees in Cartesian product networks}
\author{Rui Li, Gregory Gutin, He Zhang, Zhao Wang, Xiaoyan Zhang, Yaping Mao}
\date{Source: arXiv:2301.12933v3 (CC BY 4.0)}
\begin{document}
\maketitle

\noindent\emph{Source and licence.} Constructing edge-disjoint Steiner trees in Cartesian product networks. Version arXiv:2301.12933v3 (CC BY 4.0), distributed by arXiv under the Creative Commons Attribution 4.0 International licence, http://creativecommons.org/licenses/by/4.0/, as recorded in the arXiv licence metadata for this exact version and shown on the abstract page https://arxiv.org/abs/2301.12933v3. Full licence notice: \texttt{licenses/arxiv-2301.12933-CC-BY-4.0.txt}.\par\smallskip

\noindent\textit{Editorial scope.} Counted unit: the article's lemma lem2-1 (source0.tex lines 308--311). The article's complete original proof of it is delivered with it (source0.tex lines 312--419, one \texttt{proof} environment of 108 lines). No mathematics is written, repaired or re-derived: the statement and the proof are the article's own lines, in the article's own notation, and no retained line is substituted or re-flowed.  No further source-local context is needed: every cross-reference of every retained line resolves inside the two delivered spans.  One declared adaptation: exactly 1 ancillary strings inside the retained spans are omitted (image-inclusion commands and, where the source repeats a label, the repeated label definitions) and no other line differs from the source; the figure environments, with their placement, captions and labels, are kept, and the package records the omitted strings in fidelity receipts bound to the affected bodies.  No retained line contains a citation, so the wrapper carries no bibliography and no bibliography entry.  The retained lines contain the article's own diagram code, which the wrapper typesets with the standard tikz packages; no image file is delivered and the package declares no asset dependency.  Editorial formatting only, in addition: an AMS \texttt{amsart} preamble that restates the article's own declarations and macros used by the retained lines, namely the environments \texttt{lemma}, \texttt{fact} on separate counters, together with the standard packages those lines need.  The retained environments print in this document as Lemma 1, Fact 1 in order of appearance, whereas the article prints the same items with the numbering of its own full document.  The complete native source of the article as distributed in the arXiv e-print for this exact version is bundled unchanged as \texttt{sources/136696/source0.tex}.
\begin{lemma}\label{lem2-1}
If $|S_G|=|S_H|=k$, we can construct at least $a+b$
edge-disjoint $S$-Steiner trees in $G\Box H$.
\end{lemma}

\begin{proof}
Since $\lambda_k(H)=b$, it follows that there are $b$ edge-disjoint
$S_H$-Steiner trees, say $T_1, T_2,\ldots, T_b$. Without loss of
generality, let $T_1, T_2,\ldots, T_{b'}$ be the edge-disjoint
minimal $S_H$-Steiner trees such that $V(T_i)=S_H$, and $T_{b'+1},
T_{b'+2},$ $\ldots, T_{b}$ be the edge-disjoint minimal
$S_H$-Steiner trees such that $V(T_i)\supset S_H$. From the definition
of $b'$, we have $b'\leq \lfloor k/2\rfloor$. Since $b\geq a\geq
\lfloor k/2\rfloor$, it follows that $a\geq b'$. For each $i \
(1\leq i\leq k)$, let $T_1(u_i), T_2(u_i), \cdots, T_{b}(u_i)$ be
the edge-disjoint $S_H(u_i)$-Steiner trees in $H(u_i)$ corresponding
to $T_1, T_2,\ldots, T_{b}$ in $H$, where
\[
S_H(u_i)=\{(u_i,v_{j_t})\,|\,1\leq t\leq k\}.
\]


Since $\lambda_k(G)=a$, it follows that there are $a$ edge-disjoint
minimal $S_G$-Steiner trees, say $T_1', T_2',\ldots, T_a'$. For each
$j \ (1\leq j\leq k)$, let $T_1'(v_j), T_2'(v_j),$ $ \cdots,
T_{a}'(v_j)$ be the edge-disjoint $S_G(v_j)$-Steiner trees in
$G(v_j)$ corresponding to $T_1', T_2',\ldots, T_{a}'$ in $G$, where
\[
S_G(v_j)=\{(u_{i_s},v_j)\,|\,1\leq s\leq k\}.
\]


\begin{fact}\label{fact1}
For any $S_H(u_{i_s})$-Steiner tree $T_p(u_{i_s}) \ (1\leq s\leq k,
\ 1\leq p\leq b)$ and any $S_G(v_{j_t})$-Steiner tree $T_q'(v_{j_t})
\ (1\leq t\leq k, \ 1\leq q\leq a)$, we can find two edge-disjoint
$S$-Steiner trees in $(\bigcup_{s=1}^kH(u_{i_s}))\cup
(\bigcup_{t=1}^kG(v_{j_t}))$.
\end{fact}
\begin{proof}
Note that $T_q'(v_{j_1})\cup T_p(u_{i_2})\cup T_p(u_{i_3})\cup
\ldots \cup T_p(u_{i_k})$ and $T_p(u_{i_1})\cup T_q'(v_{j_2})\cup
T_q'(v_{j_3})\cup \ldots \cup T_q'(v_{j_k})$ are two edge-disjoint
$S$-Steiner trees in $(\bigcup_{s=1}^kH(u_{i_s}))\cup
(\bigcup_{t=1}^kG(v_{j_t}))$.
\end{proof}


From Fact~\ref{fact1}, we can find $2b'$ edge-disjoint $S$-Steiner
trees in
\[
\left(\bigcup_{p=1}^{b'}\bigcup_{s=1}^kT_p(u_{i_s})\right)\bigcup
\left(\bigcup_{q=1}^{b'}\bigcup_{t=1}^kT_q'(v_{j_t})\right).
\]


Since $b\geq a\geq b'$, from Fact~\ref{fact1}, we can also find
$2a-2b'$ edge-disjoint $S$-Steiner trees in
\[
\left(\bigcup_{p=b'+1}^{a}\bigcup_{s=1}^kT_p(u_{i_s})\right)\bigcup
\left(\bigcup_{q=b'+1}^{a}\bigcup_{t=1}^kT_q'(v_{j_t})\right).
\]
It remains to find $(a+b)-2b'-(2a-2b')=b-a$ edge-disjoint
$S$-Steiner trees in addition to the trees above. Note that we still have
edge-disjoint $S_H(u_{i_s})$-Steiner trees
$T_{a+1}(u_{i_s}),T_{a+2}(u_{i_s}),\ldots,T_{b}(u_{i_s})$ for each
$s \ (1\leq s\leq k)$. Observe that $T_{a+1}$ is a $S_H$-Steiner
tree such that $V(T_{a+1})\supset S_H$. Without loss of generality,
let $V(T_{a+1})-S_H=\{v_{j_{k+1}},v_{j_{k+2}},\ldots,v_{j_{k+h}}\}$.


\begin{figure}[tbp]
\centering

\caption{Graphs used in the proof of Lemma~\ref{lem2-1}.}
\label{fig:lemma-1}
\end{figure}


Note that there are $a$ edge-disjoint $S_G(v_{j_{k+c}})$-Steiner
trees $T_t'(v_{j_{k+c}})$ in $G(v_{j_{k+c}})$, where $1\leq t\leq a$
and $1\leq c\leq h$. Since $\lambda_k(H)=b$, it follows that there
are $b$ edge-disjoint $S_H$-Steiner trees, say $T_1, T_2,\ldots,
T_b$. Without loss of generality, let $T_1, T_2,\ldots, T_{b'}$ be
the edge-disjoint minimal $S_H$-Steiner trees such that
$V(T_i)=S_H$, and $T_{b'+1}, T_{b'+2},\ldots, T_{b}$ be the
edge-disjoint minimal $S_H$-Steiner trees such that $V(T_i)\supset
S_H$. From the definition of $b'$, we have $b'\leq \lfloor
k/2\rfloor$. Since $b\geq a\geq \lfloor k/2\rfloor$, it follows that
$a\geq b'$. If there is a vertex $v\in V(T_{b'+i}) \ (1\leq i\leq
b-b')$, then the degree of $v$ in $T_{b'+i}$ is at least $2$. Since
there are at most $k$ edges from $v$ to $S_H$, it follows that the
number of edge-disjoint $S_H$-Steiner tree containing $v$ is at most
$\lfloor k/2\rfloor$.


Then we have the following fact.

\begin{fact}\label{fact2}
For any $S_H(u_{i_s})$-Steiner tree $T_p(u_{i_s}) \ (1\leq s\leq k,
\ a+1\leq p\leq b)$ and any $S_G(v_{j_{k+c}})$-Steiner tree
$T_q'(v_{j_{k+c}}) \ (a+1\leq q\leq b, \ 1\leq c\leq h)$, we can
find one $S$-Steiner tree in
\[
\left(\bigcup_{s=1}^{k}H(u_{i_s})\right)\bigcup
\left(\bigcup_{c=1}^{h}G(v_{j_{k+c}})\right).
\]
\end{fact}

From Fact~\ref{fact2}, we can find $(b-a)$ edge-disjoint $S$-Steiner
trees, and the total number of edge-disjoint $S$-Steiner trees is
$a+b$, as desired.
\end{proof}

\end{document}
